\chapter{Záver}\label{chap:zaver}
Táto bakalárska práca si kládla za cieľ navrhnúť a implementovať modernú, bezpečnú a multiplatformovú mobilnú aplikáciu \uv{Office Hours} pre správu a rezerváciu konzultačných hodín na Fakulte informačných technológií VUT v Brne.

V rámci analýzy boli preskúmané existujúce riešenia – systém Wiki@Merlin, Diplomky FIT, Google Kalendár a Microsoft Excel. Každé z nich pokrýva iba časť požadovanej funkcionality, pričom žiadne neponúka súčasne ochranu rezervácií, mobilnú aplikáciu, kontrolu prístupu, viditeľnosť obsadených slotov, e-mailové notifikácie a možnosť online stretnutia (Tabuľka~\ref{tab:porovnanie}). Táto medzera sa stala motiváciou pre vývoj vlastného riešenia.

Na základe identifikovaných nedostatkov boli špecifikované funkčné a nefunkčné požiadavky a navrhnutý diagram prípadov použitia. Používateľské rozhranie vznikalo iteratívne – od ručných náčrtov cez prototyp vo Figme až po finálnu implementáciu, pričom celý proces bol riadený spätnou väzbou. Aplikácia bola implementovaná vo frameworku Flutter s využitím jazyka Dart a architektonického vzoru MVVM, ktorý zaručuje čisté oddelenie prezentačnej vrstvy od biznis logiky. Kľúčové funkcie zahŕňajú autentifikáciu pomocou jednorazového OTP kódu, správu miestností, blokov a slotov, optimistickú aktualizáciu roz\-hrania a nastavenia súkromia v súlade s GDPR. Aplikácia bola otestovaná manuálne aj používateľsky na emulátore a fyzických zariadeniach.

Stanovené ciele práce boli splnené. Výsledná aplikácia priamo odstraňuje nedostatky analyzovaných systémov: rezervácie sú chránené pred neoprávnenou manipuláciou, prístup k miestnostiam je riadený na úrovni e-mailových adries a domén, obsadené sloty sú pre Návštevníka viditeľné a systém môže aktívne informovať používateľov o zmenách, pokiaľ majú notifikácie povolené. Aplikácia bola navrhnutá a implementovaná ako multiplatformová. Zdrojový kód je plne pripravený na kompiláciu pre iOS, avšak publikácia pre túto platformu zostáva bezprostredným ďalším krokom, keďže počas tvorby práce nebol dostupný počítač s operačným systémom macOS nevyhnutným pre Xcode. Aplikácia bola úspešne nasadená pre platformu Android a publikovaná v obchode Google Play. Zdrojový kód je štruktúrovaný a zdokumentovaný tak, aby umožňoval ľahkú orientáciu a ďalší rozvoj budúcimi vývojármi. Možné smery ďalšieho rozvoja aplikácie sú podrobnejšie rozpracované v kapitole~\ref{chap:navrh}, v sekcii venovanej možnostiam rozšírenia. Súčasťou odovzdaných výstupov je taktiež plagát 
prezentujúci aplikáciu a krátke video demonštrujúce 
jej praktické použitie, ktoré boli vytvorené 
v súlade s požiadavkami zadania.